% ═══════════════════════════════════════════════════════════════
% MUSTERLÖSUNG: Strahlenschutz (Physik Klasse 10)
% Identisches Layout wie AB, Lösungen in ROT
% ═══════════════════════════════════════════════════════════════

\documentclass[10pt, a4paper]{article}

\usepackage[utf8]{inputenc}
\usepackage[T1]{fontenc}
\usepackage[ngerman]{babel}
\usepackage{helvet}
\renewcommand{\familydefault}{\sfdefault}

\usepackage[margin=1.5cm, top=1.8cm, bottom=1.5cm]{geometry}
\usepackage{enumitem}
\usepackage{tabularx}
\usepackage{booktabs}
\usepackage{array}
\usepackage{multicol}

\usepackage{xcolor}
\usepackage{tcolorbox}
\usepackage{amssymb}
\usepackage{fancyhdr}
\usepackage{lastpage}

\setlist{nosep, leftmargin=1.5em}

% Farben
\definecolor{physik}{HTML}{2c5aa0}
\definecolor{grau}{HTML}{888888}
\definecolor{hellgrau}{HTML}{f0f0f0}
\definecolor{fehler}{HTML}{c62828}
\definecolor{fehlerbg}{HTML}{ffebee}
\definecolor{loesung}{HTML}{c62828}

\newcommand{\fachfarbe}{physik}

% Variablen
\newcommand{\thema}{Strahlenschutz}
\newcommand{\fach}{Physik}
\newcommand{\klasse}{Klasse 10}
\newcommand{\stunde}{Stunde 9-10}

% Kopf- und Fußzeile
\pagestyle{fancy}
\fancyhf{}
\renewcommand{\headrulewidth}{0.4pt}
\lhead{\small\fach{} \klasse{} -- \thema{} \textcolor{loesung}{\textbf{[MUSTERLÖSUNG]}}}
\rhead{\small Name: \rule{4cm}{0.4pt}}
\cfoot{\small\thepage/\pageref{LastPage}}

% Befehle
\newcommand{\lsg}[1]{\textcolor{loesung}{\textbf{#1}}}

\newtcolorbox{merksatz}[1][]{
    colback=hellgrau, colframe=\fachfarbe, boxrule=0.8pt, arc=0pt,
    left=6pt, right=6pt, top=6pt, bottom=6pt,
    fonttitle=\bfseries\small, title=#1
}

\newtcolorbox{fehlerbox}{
    colback=fehlerbg, colframe=fehler, boxrule=0.8pt, arc=0pt,
    left=6pt, right=6pt, top=4pt, bottom=4pt,
    fonttitle=\bfseries\small, title=Typische Fehler
}

\newcommand{\abschnitt}[1]{%
    \vspace{8pt}\noindent\textbf{\textcolor{\fachfarbe}{#1}}\vspace{4pt}
}

\begin{document}

\begin{center}
{\large\bfseries Arbeitsblatt: \thema{} -- \stunde{} \textcolor{loesung}{[MUSTERLÖSUNG]}}\\[2pt]
{\small\textcolor{grau}{Begleitend zum Lernpfad -- in den Hefter übernehmen}}
\end{center}
\vspace{-2pt}\hrule\vspace{8pt}

\begin{multicols}{2}
\small

% ─────────────────────────────────────────────────────────────
% KASTEN 1: Dosisbegriffe
% ─────────────────────────────────────────────────────────────

\abschnitt{Kasten 1: Dosisbegriffe}

\begin{merksatz}[Einheiten der Strahlung]
\renewcommand{\arraystretch}{1.6}
\begin{tabular}{|p{2.2cm}|p{1.5cm}|p{2.3cm}|}
\hline
\textbf{Begriff} & \textbf{Einheit} & \textbf{Bedeutung} \\
\hline
Aktivität & \lsg{Bq} & \lsg{Zerfälle/s} \\
\hline
Energiedosis & \lsg{Gy} & \lsg{Energie/Masse} \\
\hline
Äquivalentdosis & \lsg{Sv} & \lsg{Biol. Wirkung} \\
\hline
\end{tabular}
\renewcommand{\arraystretch}{1}

\vspace{6pt}

\textbf{Umrechnung:}

1 Sv = \lsg{1000} mSv = \lsg{1.000.000} µSv
\end{merksatz}

\vspace{6pt}

% ─────────────────────────────────────────────────────────────
% KASTEN 2: Grenzwerte
% ─────────────────────────────────────────────────────────────

\abschnitt{Kasten 2: Grenzwerte und Dosen}

\begin{merksatz}[Übersicht]
Ergänze die Werte aus dem Lernpfad:

\vspace{4pt}
\renewcommand{\arraystretch}{1.4}
\begin{tabular}{|p{3.5cm}|c|}
\hline
\textbf{Situation} & \textbf{Dosis} \\
\hline
Natürliche Strahlung (Jahr) & \lsg{2-3 mSv} \\
\hline
Grenzwert Bevölkerung (Jahr) & \lsg{1 mSv} \\
\hline
Grenzwert beruflich (Jahr) & \lsg{20 mSv} \\
\hline
Röntgen-Thorax & \lsg{0,01-0,1 mSv} \\
\hline
CT-Untersuchung & \lsg{2-20 mSv} \\
\hline
Akute Strahlenkrankheit ab & \lsg{ca. 1 Sv} \\
\hline
\end{tabular}
\renewcommand{\arraystretch}{1}
\end{merksatz}

\vspace{6pt}

% ─────────────────────────────────────────────────────────────
% KASTEN 3: Strahlenschutzgrundsätze
% ─────────────────────────────────────────────────────────────

\abschnitt{Kasten 3: Die 3 A's des Strahlenschutzes}

\begin{merksatz}[Schutzmaßnahmen]
\textbf{1. A}bstand: Je \lsg{größer} der Abstand,

desto \lsg{geringer} die Dosis.

Gesetz: Dosis $\sim$ \lsg{1/r²}

\vspace{6pt}

\textbf{2. A}bschirmung: \lsg{Geeignetes Material}

\lsg{zwischen Quelle und Person bringen}

\vspace{6pt}

\textbf{3. A}ufenthaltszeit: Je \lsg{kürzer}

die Aufenthaltszeit, desto \lsg{geringer} die Dosis.
\end{merksatz}

\columnbreak

% ─────────────────────────────────────────────────────────────
% ÜBUNG 1: Zuordnung (AFB I)
% ─────────────────────────────────────────────────────────────

\abschnitt{Übung 1: Einheiten zuordnen}

Ordne die Einheiten den Begriffen zu:

\vspace{4pt}
\textit{Bq, Gy, Sv, mSv, µSv}

\vspace{6pt}

\begin{tabular}{|p{3.5cm}|c|}
\hline
Aktivität einer Probe & \lsg{Bq} \\
\hline
Energiedosis im Gewebe & \lsg{Gy} \\
\hline
Biologische Wirkung & \lsg{Sv} \\
\hline
Jahresdosis natürlich & \lsg{mSv} \\
\hline
\end{tabular}

\vspace{8pt}

% ─────────────────────────────────────────────────────────────
% ÜBUNG 2: Berechnungen (AFB II)
% ─────────────────────────────────────────────────────────────

\abschnitt{Übung 2: Berechnungen}

\textbf{a)} Eine Person steht 2 m von einer Quelle

entfernt und erhält 100 µSv/h.

Welche Dosis erhält sie bei 4 m Abstand?

\vspace{4pt}
\lsg{4m = 2×2m → Faktor 4 weniger → 100/4 = 25 µSv/h}

\vspace{6pt}

\textbf{b)} Ein Arbeiter darf max. 20 mSv pro Jahr

erhalten. Bei seiner Tätigkeit erhält er

0,1 mSv pro Stunde. Wie viele Stunden

darf er maximal arbeiten?

\vspace{4pt}
\lsg{20 mSv ÷ 0,1 mSv/h = 200 Stunden}

\vspace{8pt}

% ─────────────────────────────────────────────────────────────
% ÜBUNG 3: Fehler erklären (AFB III)
% ─────────────────────────────────────────────────────────────

\abschnitt{Übung 3: Fehler erklären (AFB III)}

\begin{fehlerbox}
Erkläre, warum diese Aussage falsch ist:

\vspace{4pt}

\textit{„Jede Strahlung ist schädlich, deshalb sollte man Röntgenuntersuchungen immer ablehnen."}

\vspace{4pt}
\lsg{Falsch! Bei medizinischen Untersuchungen überwiegt oft der Nutzen (Früherkennung, Diagnose) das geringe Risiko. Ärzte wägen ab (Rechtfertigung).}
\end{fehlerbox}

\vspace{8pt}

% ─────────────────────────────────────────────────────────────
% ÜBUNG 4: Transfer (AFB III)
% ─────────────────────────────────────────────────────────────

\abschnitt{Übung 4: Anwendung}

Begründe: Warum tragen Röntgenärzte eine Bleischürze, obwohl sie selbst nicht geröntgt werden?

\vspace{4pt}
\lsg{Streustrahlung! Ein Teil der Röntgenstrahlen wird vom Patienten in alle Richtungen gestreut. Bei vielen Aufnahmen pro Tag summiert sich die Dosis.}

\end{multicols}

% ═══════════════════════════════════════════════════════════════
% SEITE 2: FALLSTUDIE
% ═══════════════════════════════════════════════════════════════

\newpage

\begin{center}
{\large\bfseries Fallstudie: Strahlenschutz in der Praxis \textcolor{loesung}{[MUSTERLÖSUNG]}}\\[2pt]
{\small\textcolor{grau}{Anwendung der 3 A's}}
\end{center}
\vspace{-2pt}\hrule\vspace{12pt}

\abschnitt{Fall A: Röntgenabteilung im Krankenhaus}

\begin{merksatz}[Szenario]
Eine Röntgenassistentin macht täglich 20 Röntgenaufnahmen. Bei jeder Aufnahme erhält sie ohne Schutzmaßnahmen eine Streustrahlung von 10 µSv.
\end{merksatz}

\vspace{6pt}

\textbf{1. Berechne die Jahresdosis ohne Schutz} (250 Arbeitstage):

\vspace{4pt}
\lsg{20 Aufnahmen × 10 µSv × 250 Tage = 50.000 µSv = 50 mSv/Jahr}

\vspace{8pt}

\textbf{2. Liegt sie über dem Grenzwert?} \lsg{Ja! 50 mSv > 20 mSv (beruflicher Grenzwert)}

\vspace{8pt}

\textbf{3. Welche Schutzmaßnahmen reduzieren die Dosis? Nenne 3 konkrete Maßnahmen:}

\vspace{4pt}
\lsg{1. Abstand: Hinter Bleiglaswand stehen, nicht im Raum bleiben}

\vspace{4pt}
\lsg{2. Abschirmung: Bleischürze, Schilddrüsenschutz, Bleiglasbrille tragen}

\vspace{4pt}
\lsg{3. Aufenthaltszeit: Raum während Aufnahme verlassen, Fernauslöser verwenden}

\vspace{12pt}

\abschnitt{Fall B: Transport radioaktiver Substanzen}

\begin{merksatz}[Szenario]
Ein Krankenhaus erhält radioaktives Jod-131 ($T_{1/2}$ = 8 Tage) für Schilddrüsendiagnostik. Die Aktivität bei Lieferung beträgt 800 MBq.
\end{merksatz}

\vspace{6pt}

\textbf{1. Welche Aktivität hat das Präparat nach 16 Tagen?}

\vspace{4pt}
\lsg{16 Tage = 2 Halbwertszeiten → 800 → 400 → 200 MBq}

\vspace{8pt}

\textbf{2. Welche Sicherheitsmaßnahmen müssen beim Transport beachtet werden?}

\vspace{4pt}
\lsg{Abgeschirmte Transportbehälter (Blei), Kennzeichnung mit Strahlenwarnzeichen}

\vspace{4pt}
\lsg{Dokumentation, schneller Transport, Schulung des Personals, Dosimetrie}

\vspace{12pt}

\abschnitt{Fall C: Abstandsgesetz anwenden}

\begin{merksatz}[Szenario]
Bei einer Industrieröntgenanlage beträgt die Dosisleistung in 1 m Abstand 500 µSv/h.
\end{merksatz}

\vspace{6pt}

\textbf{Vervollständige die Tabelle} (Abstandsquadratgesetz):

\vspace{6pt}

\begin{tabular}{|c|c|c|}
\hline
\textbf{Abstand} & \textbf{Faktor} & \textbf{Dosisleistung} \\
\hline
1 m & 1 & 500 µSv/h \\
\hline
2 m & \lsg{1/4} & \lsg{125 µSv/h} \\
\hline
5 m & \lsg{1/25} & \lsg{20 µSv/h} \\
\hline
10 m & \lsg{1/100} & \lsg{5 µSv/h} \\
\hline
\end{tabular}

\vspace{12pt}

\textbf{Fazit:} In welchem Abstand ist die Dosisleistung unter 1 µSv/h?

\vspace{4pt}
\lsg{500/r² < 1 → r² > 500 → r > 22,4 m. Also ab ca. 23 m Abstand.}

\end{document}
